\section{Implementation in the BX3 Framework}\label{sec:bailout-implementation}

The Bailout Protocol, as specified in \S~\ref{sec:bailout-full}, is not an
independent mechanism operating alongside the BX3 framework.  It is the
runtime execution layer of BX3's upstream accountability guarantee ---
the mechanism that converts the static promise ``exceptions reach the human
anchor'' into a deterministic, auditable, architecturally enforced propagation
event.  This section describes how the protocol is embedded within the three-layer
model, how each layer participates in its operation, and how the forensic record
generated by the protocol constitutes the canonical accountability artifact of
the entire framework.

% ---------------------------------------------------------------
\subsection{The Purpose Layer as Human Accountability Anchor}\label{sec:bailout-purpose-layer}

The \purpose{} layer carries the human principal's intent and is the exclusive
terminus for any unresolved exception in a BX3 node.  This is not a policy
choice -- it is a structural property of the layer's function, established at
node provisioning and invariant throughout the node's operational lifetime
(Invariant~\ref{inv:human-anchor}).

When the Bailout Protocol state machine reaches the \texttt{HUMAN\_ACKNOWLEDGED}
state, the \purpose{} layer receives the escalated exception with its full
diagnostic context chain appended by the propagation algorithm
(\S~\ref{sec:escalation-algorithm}).  The layer does not receive a filtered
summary or a machine-actor interpretation -- it receives the complete record
of the exception's origin, the resolution attempts that failed at each
intermediate node, the Bounds Engine outputs, and the timeout states that
triggered escalation.  This completeness is architecturally enforced by the
\fact{} layer's pre-commit hook: no machine actor may truncate, aggregate, or
redact the diagnostic context before it reaches the \purpose{} layer.

The \purpose{} layer's role in the protocol is strictly bounded.  It does not
monitor the state machine; it does not initiate escalation; it does not evaluate
whether the trigger condition has been satisfied.  It receives exceptions that
have already passed through the full propagation chain and issues one of three
binding directives: \texttt{ACK} (continue with current parameters, returning the
node to \texttt{IDLE}), \texttt{RESOLVE} (execute a specific resolution prescribed
by the human principal, entering \texttt{RESOLVED}), or \texttt{REJECT}
(disallow the proposed action and enter a quiescent error state, also entering
\texttt{RESOLVED}).  The directive is recorded in the authorization ledger via
the \fact{} layer before it takes effect; no machine actor may issue,
simulate, or substitute a directive.

The immutable human anchor $h(n)$ is provisioned at node initialization and
stored in \fact-layer-managed memory.  No machine actor may redirect the
bailout terminus to a different principal, insert itself as an intermediate
resolver, or modify the human-root identifier.  This is the architectural
realization of the upstream accountability guarantee's terminal commitment:
the exception reaches \emph{that specific human}, no one else, with no
substitutes permitted.

% ---------------------------------------------------------------
\subsection{Bounds Engine Triggering}\label{sec:bailout-bounds-engine}

The \bounds{} layer carries bounded reasoning and constrained execution within
the provisioned operating range $R(n)$.  It is the layer that evaluates whether
proposed actions fall within the node's current authority and it is the layer
that detects when they do not.  The Bounds Engine is the computational component
of this layer: it executes resolution procedures, produces candidate resolutions,
and emits the trigger condition that activates the Bailout Protocol.

The primary trigger condition TC (Equation~\ref{TC}) fires when three conditions
hold simultaneously: the observable input $x$ falls outside $R(n)$; $x$ has not
been previously resolved in the current operational session; and the Bounds
Engine's self-assessed confidence in its ability to resolve $x$ falls below the
provisioned threshold $\tau$.  Critically, the Bounds Engine does not exercise
discretion over this evaluation -- it performs a deterministic read against
the operating range definition and the confidence function, both provisioned at
node initialization and immutable at runtime.  There is no judgment call: the
trigger condition either fires or it does not, and the result directly drives the
state machine transition without an intervening machine-actor gate.

When TC fires, the state machine transitions from \texttt{IDLE} to
\texttt{BOUNDED} (T1), and the Bounds Engine begins a time-bounded resolution
attempt with timeout $T_{\text{bounds}}$.  During the \texttt{BOUNDED} state, all
resolution paths within current authority remain open; the Bounds Engine may
attempt multiple strategies in sequence or in parallel, each submission reviewed
by the \fact{} layer against the authorization ledger.  If the \fact{} layer
approves a resolution before $T_{\text{bounds}}$ expires, the node returns to
\texttt{IDLE} via transition T2.  If all resolution paths are exhausted or
$T_{\text{bounds}}$ expires, the state machine advances to \texttt{BAIL\_PENDING}
via T3.

The secondary trigger is the explicit \texttt{bailout\_signal}, which the Bounds
Engine may emit at any time regardless of confidence scores.  This handles the
case where the Bounds Engine encounters a condition that is within $R(n)$ but
logically unresolvable -- for example, two inputs within the operating range
that contradict each other in a way that cannot be arbitrated without the
principal's intent.  The \texttt{bailout\_signal} bypasses TC entirely and
immediately initiates the transition to \texttt{ESCALATING} via T3 and T4,
ensuring that logical contradictions cannot be papered over by confidence
inflation.

The no-discretion property is architecturally enforced at the pre-commit stage:
the \fact{} layer's gate rejects any retroactive resolution attempt submitted
after $T_{\text{bounds}}$ has expired, preventing the Bounds Engine from
issuing a late resolution to suppress a warranted escalation.  This check
operates identically on all machine actors; there is no privileged actor with
the ability to bypass it.

% ---------------------------------------------------------------
\subsection{Forensic Ledger Recording}\label{sec:bailout-forensic-ledger}

The \fact{} layer maintains the Forensic Ledger as the append-only record of
all protocol events across all five pillars, and the Bailout Protocol generates
the most consequential class of entries in that record.  Every state transition
in $\mathcal{B}$, every trigger evaluation, every timeout event, every pathway
selection decision, every \texttt{SendToParent} invocation, and every human
directive received from $h(n)$ is recorded in the Forensic Ledger before any
subsequent action is taken by any layer of any node.

Each Forensic Ledger entry $f$ is a 6-tuple:
\[
f = (\text{id}(f),\ t_{\text{logical}}(f),\ \text{node}(f),\ \text{event}(f),\ \text{context}(f),\ \text{hash}(f))
\]
where the SHA-256 hash of the preceding entry ensures cryptographic chain
integrity.  The Forensic Ledger is not a conventional audit log -- it is the
authoritative source of truth for the runtime state of the Bailout Protocol's
state machine and the enforcement substrate for the framework's
accountability properties.  The pre-commit hooks for the Bailout Protocol,
Sandbox Gate, and Spatial Firewall all query the Forensic Ledger to determine
the current state of their respective state machines before committing any
state change.

The Bailout Protocol generates at minimum the following entry types in the
Forensic Ledger, each written atomically by the \fact{} layer before the
corresponding state transition is committed:

\begin{itemize}\itemsep=0pt
\item[\texttt{bailout-trigger}] Records the originating node, the exception
identifier, the trigger condition sub-components (bound signal, resolution
state, confidence value), and the priority function output $\pi(n, x)$.
\item[\texttt{bounds-resolution-attempt}] Records each resolution procedure
submitted by the Bounds Engine during the \texttt{BOUNDED} state, including
the authorization ledger entry reference and the \fact-layer approval or
rejection outcome.
\item[\texttt{bailout-signal}] Records explicit \texttt{bailout\_signal}
emissions, including the Bounds Engine's stated reason for the signal.
\item[\texttt{timeout-bounds-expired}] Records the expiration of
$T_{\text{bounds}}$ and the absence of a \fact-layer-approved resolution.
\item[\texttt{escalation-sent}] Records each invocation of
\texttt{SendToParent}, including the pathway selected, the payload hash,
the retry count, and the acknowledgment status.
\item[\texttt{pathway-exhausted}] Records the exhaustion of all retry
pathways and the terminal state reached.
\item[\texttt{human-ack}] Records the human anchor's acknowledgment of the
escalated exception, including the acknowledgment timestamp and the pathway
used.
\item[\texttt{human-directive}] Records the binding directive issued by
$h(n)$ -- one of \texttt{ACK}, \texttt{RESOLVE}, or \texttt{REJECT} -- and
the corresponding authorization ledger entry.
\end{itemize}

The Forensic Ledger entry for a Bailout Protocol escalation constitutes the
definitive accountability record for that exception.  It captures the complete
provenance of the exception from the triggering input through every
intermediate node's resolution attempt to the human principal's directive.
No machine actor may modify, delete, or suppress any entry; the layer's
function is to enforce the record rather than to interpret it.

The Forensic Ledger is the architectural bridge between the runtime
accountability properties of the Bailout Protocol (guaranteeing that every
exception reaches $h(n)$) and the retrospective auditability properties
(guaranteeing that this fact can be demonstrated to any external auditor
at any future point in time).  Together, they constitute accountable
escalation: exceptions are not merely resolved -- they are resolved in a
way that is permanently attributable to the human principal who resolved them.

% ---------------------------------------------------------------
\subsection{State Machine as Fact Layer Extension}\label{sec:bailout-state-machine-as-fact}

The Bailout Protocol state machine $\mathcal{B}$ is embedded in the \fact{}
layer of every BX3 node.  This is not an implementation detail -- it is a
structural necessity that places the state machine under the same architectural
constraints that govern all \fact-layer operations: deterministic enforcement,
immutable transition rules, pre-commit validation, and forensic auditability.

The state machine's six states -- \texttt{IDLE}, \texttt{BOUNDED},
\texttt{BAIL\_PENDING}, \texttt{ESCALATING}, \texttt{HUMAN\_ACKNOWLEDGED},
and \texttt{RESOLVED} -- map directly onto the operational phases of the
protocol's accountability guarantee:

\begin{itemize}\itemsep=0pt
\item[\texttt{IDLE}] represents normal bounded operation within $R(n)$;
\item[\texttt{BOUNDED}] represents an active exception condition being
evaluated within provisioned authority;
\item[\texttt{BAIL\_PENDING}] represents the formal point of no return --
the last architectural point at which a machine actor could suppress
escalation, which it cannot;
\item[\texttt{ESCALATING}] represents exception propagation toward $h(n)$
via the parent chain;
\item[\texttt{HUMAN\_ACKNOWLEDGED}] represents the human anchor's receipt
of the exception and the node's await state for a binding directive;
\item[\texttt{RESOLVED}] represents the terminal state after $h(n)$ issues
a directive.
\end{itemize}

The deterministic nature of the transition relation $\rightarrow$ (Table~1)
ensures that the protocol's behavior is predictable and auditable.  For every
$(q, \sigma) \in Q \times \Sigma$, at most one $q' \in Q$ satisfies
$(q, \sigma, q') \in \rightarrow$; determinism is enforced by the Bailout
Monitor's priority function $\pi$, which totally orders simultaneous trigger
conditions before any state change is committed.  If no transition is defined
for a $(q, \sigma)$ pair, the machine stalls and the Bailout Monitor records a
protocol violation in the Forensic Ledger -- an event that itself constitutes
an accountability record requiring human review.

The state invariants (INV-1 through INV-3) encode the protocol's core guarantees
as properties of the state machine:

\begin{align}
\text{INV}_1 &: \forall n \in \text{nodes},\ \text{state}(n) \in Q \setminus \{\text{ESCALATING}\} \iff \neg\text{active\_bailout}(n). \tag{BP-INV-1}
\end{align}

\begin{align}
\text{INV}_2 &: \text{state}(n) = \text{BAIL\_PENDING} \implies \text{active\_bailout}(n) \land \text{transition}(\text{BAIL\_PENDING}, e_{\text{timeout}}) = \text{ESCALATING}. \tag{BP-INV-2}
\end{align}

\begin{align}
\text{INV}_3 &: \text{state}(n) = \text{RESOLVED} \implies \exists d \in \{\text{ACK}, \text{RESOLVE}, \text{REJECT}\} \land \text{directive}(n) = d. \tag{BP-INV-3}
\end{align}

INV-2 is particularly significant: it encodes the no-machine-actor-discretion
property as a state machine invariant.  The transition from \texttt{BAIL\_PENDING}
to \texttt{ESCALATING} is forced by the state machine's deterministic transition
relation and enforced by the Bailout Monitor without machine-actor intervention.
The \fact{} layer's pre-commit hook ensures this transition cannot be intercepted
or redirected.

The timeout handling -- $T_{\text{bounds}}$, $T_{\text{prop}}$, $T_{\text{cap}}$,
and $T_{\text{human}}$ -- is similarly enforced by the \fact{} layer, not by the
\bounds{} layer or any machine actor.  This prevents a compromised machine actor
from artificially extending timeouts to suppress warranted escalation.  The
Pathway Health Registry, maintained by the Bailout Monitor, ensures that the
protocol selects functional pathways even under partial failure conditions,
and degraded pathways are excluded from selection until a health check succeeds.

The state machine's embedded position in the \fact{} layer ensures that the
Bailout Protocol's guarantees are not contingent on the correct behavior of the
\bounds{} layer or any other machine actor.  The \fact{} layer enforces the
protocol's transition relation regardless of what any machine actor attempts --
it is the structural enforcement substrate that makes the upstream accountability
guarantee architecturally compulsory rather than operationally optional.

% ---------------------------------------------------------------
\subsection{Relationship to the Upstream Accountability Guarantee}\label{sec:bailout-uag}

The BX3 upstream accountability guarantee -- that systems fail upward into human
consciousness, never downward into autonomous chaos -- is realized through the
coordinated operation of all five pillars, but the Bailout Protocol is its
runtime expression.  Loop Isolation prevents distortions in the causal chain
preceding the exception; Recursive Spawning maintains the parent-child
accountability chain that serves as the propagation substrate; Spatial Firewall
enforces domain separation during propagation; Sandbox Gate provides controlled
inter-domain communication channels; and the Bailout Protocol converts the
static guarantee into a deterministic propagation event with immutable forensic
attribution.

Without the Bailout Protocol, the upstream accountability guarantee would be
a structural property with no enforcement mechanism -- the human anchor $h(n)$
would be the designated terminus in principle, but nothing would compel
exceptions to reach it.  Machine actors could suppress, redirect, or silently
ignore exception states, and the guarantee would hold only for systems whose
machine actors chose to comply.  The Bailout Protocol converts the guarantee
from a structural aspiration into an architecturally enforced compulsion: the
state machine's deterministic transition relation, the \fact{} layer's
pre-commit hook, and the immutable parent-pointer architecture together ensure
that no unresolved exception state can exist in a BX3 node without propagating
to $h(n)$.

The Forensic Ledger entries generated by the protocol provide the retrospective
proof of this guarantee.  Any external auditor -- regulator, customer, investor,
or court -- can reconstruct the complete exception history of any node from
the cryptographically chained Forensic Ledger entries and verify that every
unresolved exception reached $h(n)$, that every human directive was recorded,
and that no machine actor suppressed, redirected, or substituted for the human
principal.  The guarantee is not merely asserted -- it is evidenced.

The Bailout Protocol therefore occupies a unique position in the BX3 framework:
it is the mechanism that makes the upstream accountability guarantee
operationally real.  It is the translation layer between the framework's
structural commitments (three-layer separation, immutable parent pointers, human
anchor designation) and the runtime execution trace (state transitions, timeout
events, propagation paths, human directives) that constitutes the definitive
accountability record of any BX3 system's operation.